\section{Кластерное разложение для Ti--Al на ГПУ-решётке}
\label{sec:ce-model}

\subsection{Что аппроксимирует модель C46}

Для заданного значения объёма на атом \(v\) химическую конфигурацию
\(\sigma\) зададим переменными \(s_i=-1\) для Ti и \(s_i=+1\) для Al
в узлах исходной ГПУ-решётки. Доля алюминия \(x\) определяет среднее
\(\langle s_i\rangle=2x-1\). Модель C46 использует выборку из 46
структур и 20 кластерных функций в формализме Санчеса, Дюкастеля и
Гратиаса~\cite{sanchez1984}. Эффективная энергия конфигурации на
атом, которую затем используют при расчёте конфигурационного равновесия, равна
\begin{equation}
 f_{\mathrm{CE}}(\sigma,T,v)
 =\sum_{\gamma}X_{\sigma\gamma}
 \bigl[J_\gamma^{\mathrm{stat}}(v)+
 J_\gamma^{\mathrm{vib}}(T,v)\bigr],
 \qquad
 X_{\sigma\gamma}=\frac{m_\gamma^{\mathrm{cell}}}{2}
 \left\langle\prod_{i\in\gamma}s_i\right\rangle_{\sigma,\gamma}.
 \label{eq:ce-model-normalization}
\end{equation}
Среднее берут по трансляциям и симметрически эквивалентным положениям
кластера. Кратность \(m_\gamma^{\mathrm{cell}}\) относится к примитивной
ячейке из двух атомов. Поэтому точечная функция равна \(2x-1\), а пустая
функция в принятой нормировке равна \(1/2\). Коэффициенты эффективных
кластерных взаимодействий (ECI) имеют единицы энергии на кластер.
Конфигурационная энтропия возникает при последующем усреднении по
заселённостям узлов. Колебательная энтропия уже входит в
\(J_\gamma^{\mathrm{vib}}\).

Базис включает пустой и точечный кластеры, семь парных оболочек
\(p_1,\ldots,p_7\) и одиннадцать треугольных орбит
\(112,114,117,123,125,136,222_1,222_2,226,247,334\).
Три цифры указывают номера оболочек трёх рёбер. Два базальных треугольника
\(222\) различаются окружением. В проекции на базальную плоскость
узел соседнего слоя лежит над центром одного из них. Порядок этих орбит в
авторском наборе и в выходном файле ATAT различается. После геометрического
сопоставления орбит расхождение по всем \(46\times20=920\) элементам
авторской матрицы \(X\) составило около \(5{,}34\times10^{-14}\). В
независимой проверке на 39 статических структурах экспорт \(J/2\) в
родном порядке ATAT воспроизвёл \(XJ\) через \texttt{corrdump -eci}
с максимальным расхождением \(1{,}2\times10^{-14}\)~эВ/атом. Такая
проверка обнаруживает как перестановку орбит, так и ошибку в пересчёте
из примитивной ячейки на атом.

\subsection{Выбор структур и подготовка обучающих энергий}

В обеих выборках исходное разбиение содержит 39 обучающих и семь контрольных
структур. Для независимой выборки ATAT перечисление ячеек с числом
узлов до десяти дало 186 Ti-богатых структур. На всех этих структурах
столбцы двух треугольников \(222\) совпадают, и ранг матрицы равен 19.
Расширение до 12 узлов дало 5777 кандидатов, из них 1174 при
\(x\leq0{,}35\). Для 944 кандидатов удалось подобрать фононную суперъячейку
в принятых ограничениях. Отобрали 40 малых структур, включая чистый
\(\alpha\)-Ti и упорядоченный \(\mathrm{Ti}_3\mathrm{Al}\) типа
\(D0_{19}\), и шесть 64-узельных конфигураций. Последние образуют три
пары с 17, 20 и 22 атомами Al, в каждой паре различается степень
химического порядка при одном составе. Для исходных 39 обучающих строк
ранг обеих матриц равен 20. Число обусловленности в нормировке
(\ref{eq:ce-model-normalization}) составляет 263 для C46 и 104 для ATAT.
В независимой выборке восемь конечных конфигураций совпали с
авторскими с учётом симметрии.

Статическую метку для независимой выборки получают релаксацией атомных
координат при заданной ячейке на том же MTP до остаточной силы
\(10^{-5}\)~эВ/\AA. Корреляции относят к исходным узлам. Из пар
смещений \(\pm0{,}01\)~\AA\ получают силовые постоянные, вычитая
силы несмещённой суперъячейки и налагая акустическое правило сумм.
Квантовая гармоническая свободная энергия на атом имеет вид
\begin{equation}
 f_{\mathrm{vib}}(T)=\frac{1}{n}
 \sum_{\mathbf q,j}'w_{\mathbf q}
 \left[\frac{h\nu_{\mathbf qj}}{2}
 +k_{\mathrm B}T\ln\!\left(1-e^{-h\nu_{\mathbf qj}/(k_{\mathrm B}T)}\right)
 \right].
 \label{eq:ce-model-phonon}
\end{equation}
Здесь \(n\) --- число атомов в исходной структуре,
\(\sum_{\mathbf q}w_{\mathbf q}=1\). Из суммы исключают три трансляции
в точке \(\Gamma\). Сетки с центром в \(\Gamma\) имели размер
\(16^3\) для структур до 32 атомов и \(6^3\) для 64-атомных структур.
Зависимость гармонических частот от объёма вводит в расчёт
квазигармоническое расширение. При заданном объёме эта модель учитывает
кривизну потенциальной энергии вблизи релаксированной структуры.

Использовали девять узлов \(v/v_0\) от 0,94 до 1,12 при
\(v_0=17{,}28833\)~\AA\({}^3\)/атом. В каждом узле статические и
колебательные метки в нашем коде аппроксимировали отдельно обычным
методом наименьших квадратов с одинаковым весом структуры и общей
маской фононной устойчивости. Авторы подгоняли статическую и полную
гармоническую энергии. Линейность регрессии при одинаковой маске
обеспечивает эквивалентность двух процедур. Из \(46\times9=414\)
сочетаний структуры и объёма в
авторском наборе допустимы 409, в независимом --- 396. При \(v_0\)
в обучение входят все 39 авторских структур и 37 независимых. Две
отброшенные независимые структуры дали мнимые внутренние частоты с
модулями 1,3155 и 0,1786~ТГц. Это существенно выше порога
\(10^{-5}\)~ТГц, принятого при обработке частот, близких к нулю.
Повторная подгонка по сохранённым авторским меткам на девяти объёмах и
89 температурах (0--2200~K, шаг 25~K) воспроизвела все
\(20\times9\times89=16\,020\) полных коэффициентов архивного набора
v16 с максимальным расхождением \(2{,}85\times10^{-14}\)~эВ/кластер.
Для независимой выборки доступны 49 температур от 600 до 1800~K с тем же
шагом. Независимое вычисление меток на восьми совпавших конфигурациях
при \(v_0\) дало среднеквадратическое расхождение \(0{,}00156\)~мэВ/атом
по статической энергии и \(0{,}00104\)~мэВ/атом по колебательной
энергии при 1000~K. Повторная подгонка проверяет переход от сохранённых
меток к коэффициентам, а восемь общих декораций позволяют проверить
согласованность самих атомистических меток. Архив v16 содержит готовые
колебательные энергии. Индивидуальные фононные силы и частоты авторов
для остальных структур по этому архиву проверить нельзя.

\begin{table}[htbp]
\centering\small
\caption{Выбранные взаимодействия при 1000~K и \(v_0\), мэВ/кластер.
Второй столбец задаёт множитель при средней корреляции.}
\label{tab:ce-model-coeff}
\begin{tabular}{lrrrrr}
\toprule
 & & \multicolumn{2}{c}{C46} & \multicolumn{2}{c}{ATAT}\\
\cmidrule(lr){3-4}\cmidrule(lr){5-6}
Орбита & \(m^{\mathrm{cell}}/2\) &
\(J^{\mathrm{stat}}\) & \(J^{\mathrm{vib}}\) &
\(J^{\mathrm{stat}}\) & \(J^{\mathrm{vib}}\)\\
\midrule
\(p_1\) & 3 & 49,16 & \(-5{,}46\) & 52,79 & \(-6{,}56\)\\
\(p_2\) & 3 & 35,59 & \(-3{,}70\) & 37,27 & \(-4{,}34\)\\
\(p_4\) & 1 & 22,60 & \(-7{,}68\) & 21,20 & \(-7{,}09\)\\
\(t_{222_1}\) & 1 & 4,11 & \(-0{,}38\) & 6,14 & \(-1{,}38\)\\
\(t_{222_2}\) & 1 & \(-8{,}83\) & 2,60 & \(-5{,}70\) & 2,89\\
\(t_{123}\) & 12 & 1,57 & \(-0{,}47\) & 2,04 & \(-0{,}55\)\\
\(t_{334}\) & 3 & \(-0{,}73\) & 0,21 & \(-2{,}46\) & 0,78\\
\bottomrule
\end{tabular}
\end{table}

Положительные \(J_{p_1}\) и \(J_{p_2}\) понижают энергию разноимённых
пар. Колебательный вклад C46 при 1000~K уменьшает полное \(J_{p_1}\)
с 49,16 до 43,70~мэВ/кластер и \(J_{p_2}\) с 35,59 до 31,89.
Полные взаимодействия двух базальных треугольников противоположны
по знаку в обеих параметризациях. Для оценки влияния
взаимодействия на упорядочение важны и его кратность, и корреляции. Коэффициент
\(t_{123}\) умножается на двенадцать, а его реальный вклад определяется
изменением корреляции при упорядочении (табл.~\ref{tab:ce-model-coeff}).

\subsection{Переносимость и связь ошибки энергии с ошибкой границы}

\begin{table}[htbp]
\centering\small
\caption{Перекрёстная проверка на семи контрольных структурах каждой
выборки при 1000~K, \(v_0\). Указана среднеквадратическая ошибка
на структуру в мэВ/атом. Столбец суммы рассчитан по полным энергиям.}
\label{tab:ce-model-transfer}
\begin{tabular}{llrrr}
\toprule
Обучение & Контроль & Статика & Колебания & Сумма\\
\midrule
C46  & Авторский & 2,648 & 0,473 & 2,446\\
ATAT & Авторский & 2,465 & 0,700 & 2,121\\
C46  & ATAT      & 2,231 & 0,300 & 2,260\\
ATAT & ATAT      & 2,483 & 0,837 & 2,724\\
\bottomrule
\end{tabular}
\end{table}

При \(v_0\) среднеквадратическая ошибка статической подгонки равна
2,207~мэВ/атом для C46 и 3,354~мэВ/атом для ATAT на общей для двух
компонент маске каждой выборки. При последовательном исключении одной
обучающей структуры (LOOCV) эти ошибки возрастают до 5,117 и
9,200~мэВ/атом. Следовательно, улучшение числа обусловленности матрицы
само по себе не гарантирует переносимость усечённого базиса. На обеих
контрольных семёрках ошибка полной энергии лежит в диапазоне
2,12--2,73~мэВ/атом, а порядок точности двух моделей зависит от
контрольной выборки (табл.~\ref{tab:ce-model-transfer}). Ошибки
статической и колебательной энергии одной и той же структуры могут
компенсироваться.

Пять авторских контрольных структур участвовали в выборе базиса. На
двух полностью отложенных структурах с \(22/64\) Al ошибка полной
энергии равна 1,031 для C46 и 1,047~мэВ/атом для ATAT. Среди семи
контрольных структур ATAT одна присутствует в авторском обучении.
После её исключения C46 даёт по оставшимся шести 2,359~мэВ/атом.
Разность энергий пары при одном составе чувствительнее к описанию
химического порядка. Для двух независимых контрольных конфигураций
с \(22/64\) Al гармоническая модель ATAT при 1000~K завышает разность
полных энергий между разупорядоченным и \(D0_{19}\)-подобным состояниями
на 3,631~мэВ/атом. Положение общей касательной зависит
от разности свободных энергий фаз и их производных по составу.
Поэтому воспроизведение авторских ECI до машинной точности устанавливает
воспроизводимость расчёта в данном базисе, а сравнение с прямой
атомистической термодинамикой проверяет физическую цену его предпосылок.
